\section{Introduction}
\label{sec:intro}

A coding agent does not issue one request. It issues a \emph{session}: a multi-round job
in which every round ends with a tool call, and every tool call returns with more context
than the round before. The session is long-lived, its state is expensive to rebuild, and
it is reused. On a serving engine that was designed for independent requests, none of
these three properties is visible.

We measured this on a replay of 19 long coding-agent sessions (602 rounds, median 28
rounds per session, longest 81) collected from Claude and Codex conversations. A round
arrives with a median 48.5K-token context (p90 71.1K, max 92.8K). The terminal working
set---what all 19 sessions would occupy if every one of them were live---is 1.31M tokens.
The KV pool of our engine, a 7B model served with a 96K context on a single 24\,GB GPU,
holds 101.5K tokens. The working set is therefore $12.9\times$ the pool, and the average
session alone needs 69K tokens: \emph{the pool holds 1.47 sessions}.

The cost of ignoring this is paid every round. A round's new content is a median of 990
tokens, but it arrives attached to a 44.9K-token prefix. Serving each round as an
independent request re-reads that prefix: across the trace, the engine is asked to
process $13.0\times$ more prefix than it appends, and a system with no reuse at all would
do $14.0\times$ the prefill work of a perfect one. The trace makes the same point in the
other direction: it is cheap to keep a session, and expensive to rebuild it.

The obvious response---evict and recompute---runs into the shape of the workload rather
than its size. Tool waits are bimodal: the median gap between rounds is 181\,ms, the mean
is 15.9\,s, and the 13.1\% of gaps longer than ten seconds account for 93.4\% of all
parked time. A session's idle period therefore tells you almost nothing about the next
one: half of the time the agent is back before an eviction decision could pay for itself,
and the other half it is gone for tens of seconds. A policy that suspends on every gap
recomputes constantly; a policy that never suspends holds sessions it cannot afford.

\paragraph{Prior fixes treat this as placement, and placement forces prediction}
Paged attention, prefix caching, and tiered KV stores decide \emph{where} a block lives;
their admission counterpart is a replacement decision that presupposes a guess about the
future---which prefix will be reused, when a session will return, how long its next round
will be. The systems that address agentic workloads keep that presupposition: TTL-based
scheduling predicts a session's reuse window, program-aware admission predicts the next
step's demand, feedback controllers adjust a concurrency limit after congestion is
observed, and overload handlers shed sessions outright. Each is a threshold on a
predicted quantity. None of them states which set of sessions is \emph{feasible} under
the capacity it actually has, and none prices what it costs to bring a suspended session
back.

\paragraph{Our position: make the session suspendable, then stop predicting}
If a session can be paused at a round boundary and resumed by \emph{copying} its KV back
from a host tier, then the question ``should this session be admitted?'' stops being a
forecast and becomes an accounting check against a ledger of bytes:
\begin{quote}
\emph{Admission is a feasibility test, not a threshold.}
\end{quote}
Three things follow, and they organise the paper. First, feasibility is stated in
quantities we can measure now---the session's projected footprint, the capacity of each
tier, the bandwidth between them---so no predictor is on the critical path. Second,
suspension is not free and not always right: reviving a session costs a copy, recomputing
it costs a prefill, and the two are comparable for short contexts, which yields a
per-session break-even instead of a policy constant. Third, the one genuinely uncertain
quantity, how much a session will grow before it returns, is small enough---a few percent
of the footprint it protects---that a crude estimate of it cannot move the ledger, which is
why the design needs no predictor.

\paragraph{Contributions}
\begin{enumerate}
\item \textbf{Formulation} (\S\ref{sec:model}): we cast admission for suspendable agent
sessions as a feasibility problem over a two-tier capacity ledger, with an explicit
revival cost. The formulation replaces a threshold on predicted load with a test on
measured capacity, and it exposes the one parameter that remains---the fraction of the
pool held in reserve (\S\ref{sec:design}).
\item \textbf{Mechanism} (\S\ref{sec:design}--\ref{sec:impl}): named permits make the
ledger enforceable at request granularity; reversible suspension at round boundaries
keeps a session's state in the host tier instead of discarding it; a growth reserve
estimated from the session's own history replaces a point forecast; and the admission
accounting is dominated by measured bytes rather than by that estimate. We implement the design on SGLang with a
hierarchical KV cache and a 96K-context 7B model.
\item \textbf{Analysis and evidence} (\S\ref{sec:analysis}, \S\ref{sec:eval}): we give a
feasible-set bound, a session-level revival-versus-recompute break-even, and a bound on
the projection error; we evaluate request-level serving, an agent-level
feedback controller, and our design on the same trace, engine, and memory budget; and we
report the regimes in which suspension does not pay.
\end{enumerate}

\paragraph{Results in one line}
\pend{one-sentence headline with the three arms, e.g.\ ``on the 602-round replay our
design finishes the batch in X minutes against Y (request-level) and Z (agent-level),
while meeting the per-round deadline on W\% of rounds against V\% and U\%''}.

\paragraph{Scope}
Our testbed is a single 24\,GB GPU running a 7B model with a 96K context. We choose this
point because it makes the regime sharp, not because the mechanism is specific to the
card: the binding variable is the ratio of the live working set to KV capacity, and when
that ratio is small, admission is not the problem. We do not claim that our constants
transfer to larger devices or to clusters; we do claim that the ledger, the break-even,
and the projection bound are stated in quantities---bytes, bandwidth, session
horizon---that do transfer (\S\ref{sec:discussion}).
